\section*{第 \ref{ch:b2:series} 章——数列与级数}

\begin{solution}{exo:b2:series:1}
$\sum\frac{\cos n}{n}$：应用阿贝尔判别法，取 $a_n = \frac1n$，$b_n =
\cos n$，其部分和有界（是几何和的实部，
见\cref{ex:b2:series:sinn}）：收敛（不\omterm{def:b2:series:def}{绝对收敛}，
由同样的 $\cos^2$ 技巧）。

$\sum \frac{(-1)^n}{\ln n}$（$n \geq 2$）：交错级数判别法，
$\frac{1}{\ln n} \downarrow 0$：收敛；不\omterm{def:b2:series:def}{绝对收敛}（$\ln n
\leq n$）。

$\sum \frac{(-1)^n}{n^{3/4} + \cos n}$：展开，
\[
\frac{(-1)^n}{n^{3/4} + \cos n}
= \frac{(-1)^n}{n^{3/4}}\cdot
\frac{1}{1 + \frac{\cos n}{n^{3/4}}}
= \frac{(-1)^n}{n^{3/4}}
- \frac{(-1)^n\cos n}{n^{3/2}} + O\Bigl(\frac{1}{n^{9/4}}\Bigr).
\]
第一个级数：交错级数，收敛。第二个：\omterm{def:b2:series:def}{绝对
收敛}（量级 $\frac{1}{n^{3/2}}$）。第三个：\omterm{def:b2:series:def}{绝对
收敛}。于是原级数收敛。
\end{solution}

\begin{solution}{exo:b2:series:2}
$\sum_{m\geq0} z^m$ 与自身的\omterm{thm:b2:series:fubini}{柯西乘积}（当 $\abs z < 1$ 时二者均
\omterm{def:b2:series:def}{绝对收敛}）：对角线系数为 $c_k =
\sum_{m=0}^{k} 1 = k + 1$，于是
\[
\frac{1}{(1-z)^2} = \sum_{k\geq0} (k+1)z^k .
\]
乘以 $z$ 并换标号，得 $\sum_{n \geq 1} n z^n =
\frac{z}{(1-z)^2}$。
\end{solution}

\begin{solution}{exo:b2:series:3}
设 $B_n = \sum_{k \leq n} b_k \to B$。取 $a_k = k$，用\omterm{thm:b2:series:abel}{阿贝尔分部求和}：
\[
\sum_{k=1}^{n} k\,b_k = n B_n - \sum_{k=1}^{n-1} B_k
\quad\Longrightarrow\quad
\frac1n \sum_{k=1}^{n} k b_k = B_n - \frac{1}{n}\sum_{k=1}^{n-1}
B_k .
\]
收敛数列 $(B_k)$ 的切萨罗平均收敛于其极限 $B$
（参见本科一年级分册），故右端趋于 $B - B = 0$。
\end{solution}

\begin{solution}{exo:b2:series:4}
若 $\theta \in \pi\Z$：各项均为零——显然收敛。
设 $\theta \notin \pi\Z$。

$\alpha > 1$：\omterm{def:b2:series:def}{绝对收敛}（以 $n^{-\alpha}$ 为优级数）。

$0 < \alpha \leq 1$：可用阿贝尔判别法（取 $a_n = n^{-\alpha}
\downarrow 0$，$\sin n\theta$ 的部分和以
$\frac{1}{\abs{\sin(\theta/2)}}$ 为界，由几何和得）：收敛。不
\omterm{def:b2:series:def}{绝对收敛}：$\abs{\sin n\theta} \geq \sin^2 n\theta = \frac{1 -
\cos 2n\theta}{2}$，而 $\sum \frac{1 - \cos 2n\theta}{2n^\alpha}$
发散（$\sum n^{-\alpha}$ 发散；$\sum
\frac{\cos 2n\theta}{n^\alpha}$ 由阿贝尔判别法知收敛（当 $2\theta
\notin 2\pi\Z$ 时）；例外情形 $2\theta \in 2\pi\Z$ 即
$\theta \in \pi\Z$，已于前面处理）。条件收敛。
\end{solution}

\begin{solution}{exo:b2:series:5}
设 $p_1, p_2, \dots$ 为 $(u_n)$ 中依次出现的非负项，
$q_1, q_2, \dots$ 为其中的负项。$\sum p_k$ 与 $\sum
q_k$ 均发散：若其中一个收敛，则另一个
等于收敛级数 $\sum u_n$ 减去前者，从而也收敛——于是
$\sum \abs{u_n} = \sum p_k - \sum q_k$ 将收敛，与
假设矛盾。又 $u_n \to 0$（因 $\sum u_n$ 收敛）。

贪心重排：依次取正项 $p_1, p_2, \dots$，直到
部分和首次超过 $\ell$（可行，因 $\sum p_k =
+\infty$）；随后取负项，直到部分和首次降至
$\ell$ 以下（可行，因 $\sum q_k = -\infty$）；无限重复此过程（每
一阶段均为有限步，且每一项恰用一次：确为真正的
重排）。每次换向后，部分和与
$\ell$ 的距离至多为最后所用的项。由于第 $m$ 次换向时所用项的
指标 $\to \infty$，且 $u_n \to 0$，
部分和收敛于 $\ell$。
\end{solution}

\begin{solution}{exo:b2:series:6}
可和性：$\frac{\abs x^{m+n}}{m!n!}$ 的有限部分和
以 $\bigl(\sum_m
\frac{\abs x^m}{m!}\bigr)^2 = \eu^{2\abs x}$ 为界。按对角线分组
（见\cref{thm:b2:series:fubini}）：
\[
(\eu^{x})^2 = \sum_{m,n} \frac{x^{m+n}}{m!\,n!}
= \sum_{k=0}^{\infty} x^k \sum_{m+n=k} \frac{1}{m!\,n!}
= \sum_k \frac{x^k}{k!}\sum_{m=0}^{k}\binom km
= \sum_k \frac{(2x)^k}{k!} = \eu^{2x} .
\]
\end{solution}

\begin{solution}{exo:b2:series:7}
可和性：作为乘积族，以 $\zeta(2)^2$ 为界
（见\cref{thm:b2:series:fubini}中的\omterm{thm:b2:series:fubini}{柯西乘积}论证）。映射
$(q, a, b) \mapsto (qa, qb)$，由满足 $\gcd(a, b) = 1$ 的三元组
到数对 $(m, n)$ 的映射是双射（取 $q = \gcd(m,n)$）。按此对
\omterm{def:b2:series:summable}{可和族}分组（即指标集的一个划分——
对\omterm{def:b2:series:summable}{可和族}而言是允许的，依据
\cref{thm:b2:series:rearrangement}/\cref{thm:b2:series:fubini}
应用于可数个类的划分）：
\[
\zeta(2)^2 = \sum_{q} \sum_{\gcd(a,b)=1} \frac{1}{q^4 a^2 b^2}
= \zeta(4)\, S,
\qquad\text{于是}\qquad
S = \frac{\zeta(2)^2}{\zeta(4)} .
\]
（取 $\zeta(2) = \frac{\pi^2}{6}$，
$\zeta(4) = \frac{\pi^4}{90}$，由
\cref{ch:b2:fourier}可得 $S = \frac{5}{2}$。）
\end{solution}

\begin{solution}{exo:b2:series:8}
设 $A = \sum a_n$（\omterm{def:b2:series:def}{绝对收敛}），$B_m = \sum_{k\leq m} b_k \to B$，
且以 $M$ 为界。则
\[
C_N = \sum_{k=0}^{N} c_k = \sum_{n=0}^{N} a_n B_{N-n}
\]
（按 $a$ 的下标合并）。写作
\[
C_N - AB = \sum_{n=0}^{N} a_n (B_{N-n} - B) - B\sum_{n > N} a_n .
\]
末项趋于 $0$。将和式在 $n = \lfloor N/2
\rfloor$ 处分段：对 $n \leq N/2$，有 $N - n \geq N/2$，故 $\abs{B_{N-n} - B}
\leq \varepsilon_N := \sup_{m \geq N/2}\abs{B_m - B} \to 0$，且
此部分 $\leq \varepsilon_N \sum\abs{a_n}$；对 $n > N/2$，
$\abs{B_{N-n} - B} \leq 2M$，而此部分 $\leq 2M \sum_{n >
N/2} \abs{a_n} \to 0$。因此 $C_N \to AB$。
\end{solution}

\begin{solution}{exo:b2:series:9}
取 $u_n = 2^{-n} e_n$（$e_n$ 为仅有一个
$1$ 位于第 $n$ 个位置、其余为零的数列）。则 $\sum \norm{u_n}_\infty = \sum 2^{-n} <
\infty$：\omterm{def:b2:series:def}{绝对收敛}。但部分和 $S_N =
(1, \tfrac12, \dots, 2^{-N}, 0, \dots)$ 必须收敛到
数列 $(2^{-n})_n$，而它\emph{并非}最终为零：
其极限在 $E$ 之外。在 $E$ 之内，$(S_N)$ 为柯西列而无极限（$\norm{S_N
- x}_\infty \geq 2^{-N-1}$ 未能排除任何最终为零的 $x$：
对任一 $x \in E$，若其在第 $K$ 项之后全为零，则 $\norm{S_N - x} \geq
2^{-K-1}$，当 $N > K$）。故该级数在 $E$ 中不收敛。
\omterm{def:b2:metric:complete}{完备}性恰是 \cref{thm:b2:nvs:absoluteconvergence}
所需要的。
\end{solution}

\begin{solution}{exo:b2:series:10}
两者 $\arctan(n+1) - \arctan n$ 与
$\arctan\frac{1}{n^2+n+1}$ 均属于 $\intoo{0}{\frac\pi2}$，
由正切加法公式可得
\[
\tan\bigl(\arctan(n{+}1) - \arctan n\bigr)
= \frac{(n+1) - n}{1 + n(n+1)} = \frac{1}{n^2 + n + 1} :
\]
两者正切值相等，而 $\tan$ 在该区间上为单射，故
该式成立。由裂项相消，
\[
\sum_{n=1}^{N}\arctan\frac{1}{n^2+n+1}
= \arctan(N{+}1) - \arctan 1 \xrightarrow[N\to\infty]{}
\frac\pi2 - \frac\pi4 = \frac\pi4 .
\]

\end{solution}

\begin{solution}{exo:b2:series:11}
其一：$\sqrt{n+1} - \sqrt n = \frac{1}{\sqrt{n+1} + \sqrt n}
\sim \frac{1}{2\sqrt n}$，故通项 $\sim
2^{-\alpha}n^{-\alpha/2}$：收敛当且仅当 $\frac\alpha2 > 1$，
即 $\alpha > 2$。其二：$1 - \cos\frac1n \sim
\frac{1}{2n^2}$：收敛。其三：$\bigl(1 + \frac1n\bigr)^n =
\eu^{\,n\ln(1 + 1/n)} = \eu^{\,1 - \frac1{2n} + O(n^{-2})} =
\eu\bigl(1 - \frac{1}{2n} + O(n^{-2})\bigr)$，于是
\[
\eu - \Bigl(1 + \frac1n\Bigr)^{\!n} \sim \frac{\eu}{2n} :
\]
正项且与调和级数的某倍数等价，故发散。
\end{solution}

\begin{solution}{exo:b2:series:12}
由单调性，$n\,a_{2n} \leq a_{n+1} + a_{n+2} + \dots +
a_{2n} = S_{2n} - S_n \to 0$（由收敛级数的
柯西准则）。于是 $2n\,a_{2n} \to 0$，且 $(2n{+}1)\,
a_{2n+1} \leq (2n{+}1)a_{2n} = \frac{2n+1}{2n}\,(2n\,a_{2n}) \to
0$：$(na_n)$ 的两个子列均趋于 $0$，故 $na_n \to 0$。

\emph{逆命题不成立：}$a_n = \frac1{n\ln n}$ 为正的递减数列，
$na_n = \frac1{\ln n} \to 0$，而 $\sum a_n$
发散（伯特兰分界，见 \cref{pb:b2:series:1}，
第 18 题）。\emph{单调性必不可少：}令 $a_n = \frac1n$，若 $n$
为 $2$ 的幂，否则 $a_n = 2^{-n}$：$\sum a_n \leq
\sum_k 2^{-k} + \sum_n 2^{-n} < \infty$，但 $na_n = 1$
在 $2$ 的各次幂上，故 $na_n \not\to 0$。
\end{solution}

\begin{solution}{pb:b2:series:1}
\textbf{1.}对 $m$ 用归纳法：$m = 0$ 时两边均为 $1$；
归纳步乘以 $\cos\theta + \iu\sin\theta$，并用到
加法公式 $\cos(m\theta + \theta) =
\cos m\theta\cos\theta - \sin m\theta\sin\theta$，$\sin(m\theta
+ \theta) = \sin m\theta\cos\theta + \cos m\theta\sin\theta$。
改用二项式定理展开，取 $m = 2n+1$，
并合并虚部（即 $\iu\sin\theta$ 的奇次幂，
其中 $\iu^{2j+1} = (-1)^j\iu$）：
\[
\sin\bigl((2n{+}1)\theta\bigr) = \sum_{j=0}^{n}(-1)^j
\binom{2n+1}{2j+1}\cos^{2(n-j)}\theta\,\sin^{2j+1}\theta .
\]

\textbf{2.}在 $\intoo0{\frac\pi2}$ 上 $\sin\theta \neq 0$，于是各项提出
$\sin^{2n+1}\theta$，余下
$\bigl(\frac{\cos^2\theta}{\sin^2\theta}\bigr)^{n-j} =
(\cot^2\theta)^{n-j}$：即上式，其中 $P_n(x) =
\sum_j(-1)^j\binom{2n+1}{2j+1}x^{n-j}$。其 $n$ 次项
系数为 $j = 0$ 时的 $\binom{2n+1}{1} = 2n + 1 \neq 0$。

\textbf{3.}当 $\theta_k = \frac{k\pi}{2n+1}$ 时：
$\sin\bigl((2n{+}1)\theta_k\bigr) = \sin k\pi = 0$，而
$\sin^{2n+1}\theta_k \neq 0$，故 $P_n(\cot^2\theta_k) = 0$。
$\theta_k$ 在 $\intoo0{\frac\pi2}$ 中严格递增，而该区间上
$\cot^2$ 严格递减，故数值 $x_k =
\cot^2\theta_k$ 两两不同——$n$ 个互异的
根，是 $n$ 次多项式的全部根。

\textbf{4.}由韦达定理，根之和等于
$x^{n-1}$ 项与 $x^n$ 项系数之比的相反数：
\[
\sum_{k=1}^{n}\cot^2\theta_k =
\frac{\binom{2n+1}{3}}{\binom{2n+1}{1}}
= \frac{(2n+1)(2n)(2n-1)/6}{2n+1} = \frac{n(2n-1)}{3}.
\]

\textbf{5.} $\frac{1}{\sin^2\theta} = 1 + \cot^2\theta$：
求和，$n + \frac{n(2n-1)}3 = \frac{3n + 2n^2 - n}{3} =
\frac{2n(n+1)}{3}$。

\textbf{6.} 在 $\intoo{0}{\frac\pi2}$ 上：$\sin\theta < \theta <
\tan\theta$（本科一年级分册）。取倒数后不等号反向：
$\cot\theta < \frac1\theta < \frac1{\sin\theta}$，两边平方
（各项均为正）得 $\cot^2\theta < \frac1{\theta^2} <
\frac1{\sin^2\theta}$。

\textbf{7.} 在第 6 题中取 $\theta = \theta_k$，对 $k \leq
n$ 求和，利用第 4、5 题，并注意到 $\frac1{\theta_k^2} =
\frac{(2n+1)^2}{k^2\pi^2}$：
\[
\frac{n(2n-1)}3 < \frac{(2n+1)^2}{\pi^2}\sum_{k=1}^n\frac1{k^2}
< \frac{2n(n+1)}3 .
\]

\textbf{8.} 两边乘以 $\frac{\pi^2}{(2n+1)^2}$：
\[
\frac{\pi^2}{3}\cdot\frac{n(2n-1)}{(2n+1)^2}
< \sum_{k=1}^{n}\frac1{k^2} <
\frac{\pi^2}{3}\cdot\frac{2n(n+1)}{(2n+1)^2}.
\]
两个界都趋于 $\frac{\pi^2}3\cdot\frac12 = \frac{\pi^2}6$
（有理分式趋于 $\frac12$）。部分和
单调递增，故收敛，由夹逼原理得 $\zeta(2) =
\frac{\pi^2}6$：即欧拉定理，此为柯西的证法。

\textbf{9.} 部分和单调递增并趋于 $\zeta(2) =
\frac{\pi^2}6$，故 $0 \leq \frac{\pi^2}6 - \sum_{k\leq n}k^{-2}$；
由第 8 题的下界可得
\[
\frac{\pi^2}6 - \sum_{k\leq n}\frac1{k^2}
\leq \frac{\pi^2}6 - \frac{\pi^2}3\cdot\frac{n(2n-1)}{(2n+1)^2}
= \frac{\pi^2}6\cdot\frac{(2n+1)^2 - (4n^2 -
2n)}{(2n+1)^2}
= \frac{\pi^2}6\cdot\frac{6n + 1}{(2n+1)^2}
= O\Bigl(\frac1n\Bigr),
\]
与精确余项 $\frac1n - \frac1{2n^2} + O(n^{-3})$ 相符，
见\cref{exo:b2:comparison:11}。

\textbf{10.} 根的第二个初等对称函数为
$\sigma_2 = \frac{\binom{2n+1}5}{\binom{2n+1}1} =
\frac{(2n)(2n-1)(2n-2)(2n-3)}{120}$，故
\[
\sum_k\cot^4\theta_k = \sigma_1^2 - 2\sigma_2
= \Bigl(\frac{n(2n-1)}3\Bigr)^2 -
\frac{2n(2n-1)(2n-2)(2n-3)}{60}
\sim \frac{4n^4}9 - \frac{4n^4}{15} = \frac{8n^4}{45}.
\]
由夹逼 $\cot^4\theta < \theta^{-4} < (1 + \cot^2\theta)^2 =
1 + 2\cot^2\theta + \cot^4\theta$，并求和：两端外侧的和均为
$\frac{8n^4}{45}(1 + o(1))$（所加项 $n + 2\sigma_1 =
O(n^2)$ 可忽略），中间项为
$\frac{(2n+1)^4}{\pi^4}\sum_{k\leq n}k^{-4}$。由此
\[
\sum_{k\leq n}\frac1{k^4} \longrightarrow
\pi^4\cdot\frac{8/45}{16} = \frac{\pi^4}{90}.
\]

\textbf{11.} 按奇偶拆分 $\zeta(2)$：$\sum_{\text{偶}}
= \sum_j\frac{1}{(2j)^2} = \frac14\zeta(2) = \frac{\pi^2}{24}$，
故 $\sum_{\text{奇}} = \zeta(2) - \frac{\pi^2}{24} =
\frac{\pi^2}8$。交错级数：$\sum_k\frac{(-1)^{k-1}}{k^2} =
\sum_{\text{奇}} - \sum_{\text{偶}} = \frac{\pi^2}8 -
\frac{\pi^2}{24} = \frac{\pi^2}{12}$（\omterm{def:b2:series:def}{绝对收敛}
保证了重排的合法性，
见\cref{thm:b2:series:rearrangement}）。

\textbf{12.} $S = \frac{\zeta(2)^2}{\zeta(4)} =
\frac{(\pi^2/6)^2}{\pi^4/90} = \frac{90}{36} = \frac52$。
启发式地看，恒等式 $\zeta(2)^2 = \zeta(4)S$
（见\cref{exo:b2:series:7}）表明，提取 $\gcd$ 后，数对化为互素数对；
其倒数
$\frac1{\zeta(2)} = \frac6{\pi^2} \approx 0.608$ 是自然的
候选，即全部数对中互素数对的密度——
这是关于 $\lim_N \frac{1}{N^2}\#\{(m,n) \leq N :
\gcd = 1\}$ 的一个论断，其严格证明（含误差项）属于
本科三年级分册。

\textbf{13.} 直接求和的误差为 $\sim \frac1n$：六位
数字约需 $10^6$ 项。修正后的和
$\sum_{k\leq n}k^{-2} + \frac1n - \frac1{2n^2}$ 的误差为
$O(n^{-3})$：取 $n = 100$ 时其值为 $1.6449339\dots$，与
$\frac{\pi^2}6 = 1.6449341\dots$ 相比——误差为 $1.7\cdot10^{-7}$，
仅用百项即得七位数字。渐近修正胜过
单纯增加项数，效率高出四个数量级。

\textbf{14.} $n = 1$：$P_1(x) = 3x - 1$，根为 $\frac13$，且
确有 $\cot^2\frac\pi3 = \bigl(\frac1{\sqrt3}\bigr)^2 =
\frac13 = \frac{1\cdot1}3$。$n = 2$ 时，公式预测
$\frac{2\cdot3}3 = 2$；取 $\cos\frac\pi5 = \frac{1 +
\sqrt5}{4}$，计算得 $\cot^2 36^\circ \approx 1.894$ 及
$\cot^2 72^\circ \approx 0.106$：和为 $2.000$。

\textbf{15.} 数 $\tan^2\theta_k = \frac1{x_k}$ 是
多项式 $Q(x) = x^nP_n\bigl(\frac1x\bigr) =
\sum_{j=0}^n(-1)^j\binom{2n+1}{2j+1}x^j$ 的根（$x_k$ 均
非零）。对 $Q$ 用韦达定理：首项系数为 $(-1)^n$
（即 $j = n$ 项），次项系数为 $(-1)^{n-1}\binom{2n+1}{2n-1} =
(-1)^{n-1}\binom{2n+1}{2}$，故
\[
\sum_{k=1}^n\tan^2\theta_k =
-\frac{(-1)^{n-1}\binom{2n+1}2}{(-1)^n} = \binom{2n+1}2\cdot
\frac{2}{2n+1}\cdot\frac{2n+1}{2}
= n(2n+1).
\]
检验 $n = 1$：$\tan^2\frac\pi3 = 3 = 1\cdot3$。

\textbf{16.} 设 $\gamma = \frac{1+\alpha}2 \in
\intoo{1}{\alpha}$。于是 $\frac{n^{-\alpha}(\ln
n)^{-\beta}}{n^{-\gamma}} = n^{\gamma - \alpha}(\ln n)^{-\beta}
\to 0$（$n$ 的负幂次胜过 $\ln n$ 的任一次幂），故
从某项起，通项 $\leq n^{-\gamma}$，其中 $\gamma > 1$：
与黎曼级数比较即得收敛。

\textbf{17.} 设 $\gamma = \frac{1+\alpha}2 \in
\intoo{\alpha}{1}$。此时 $\frac{n^{-\gamma}}{n^{-\alpha}(\ln
n)^{-\beta}} = n^{\alpha-\gamma}(\ln n)^{\beta} \to 0$，故
从某项起，通项 $\geq n^{-\gamma}$，其中 $\gamma < 1$：
发散。

\textbf{18.} $f(t) = \frac1{t(\ln t)^\beta}$ 为正、
\omterm{def:b2:metric:continuity}{连续}，且当 $t$ 充分大时递减（其对数的
导数 $-\frac1t\bigl(1 + \frac{\beta}{\ln t}\bigr) < 0$
当充分大时成立）。原函数：当 $\beta \neq 1$ 时，
$\int^x f = \frac{(\ln x)^{1-\beta}}{1-\beta} + \text{常数}$，
它有有限极限，当且仅当 $\beta > 1$；当 $\beta = 1$ 时，
$\int^x f = \ln\ln x \to \infty$。
由\cref{thm:b2:comparison:seriesintegral}，级数与
积分的敛散性相同：收敛当且仅当 $\beta > 1$。

\textbf{19.} 同一判别法：$\frac{\dd}{\dd t}\ln\ln\ln t =
\frac{1}{t\ln t\,\ln\ln t}$，且 $\ln\ln\ln t \to \infty$：
发散。又 $\frac{\dd}{\dd t}\Bigl(-\frac1{\ln\ln t}\Bigr)
= \frac{1}{t\ln t\,(\ln\ln t)^2}$，且 $-\frac1{\ln\ln t} \to
0$：收敛。

\textbf{20.} 其一：$n^{1/\ln n} = \eu^{\ln n/\ln n} = \eu$，故
通项恰为 $\frac{1}{\eu\,n}$，即调和级数的
倍数，\emph{发散}——指数 $1 +
\frac1{\ln n}$ 趋于 $1$ 太快。其二：$n^{1/\ln\ln n}
= \eu^{\ln n/\ln\ln n}$，且 $\frac{\ln n}{\ln\ln n} \geq
2\ln\ln n$，从某项起，于是 $n^{1/\ln\ln n} \geq (\ln n)^2$：
通项 $\leq \frac1{n(\ln n)^2}$，即收敛的伯特兰
级数（第 18 题）：\emph{收敛}。分界线
严格位于这两个指数之间。

\textbf{21.} $R_n \downarrow 0$，且
\[
\sqrt{R_n} - \sqrt{R_{n+1}} = \frac{R_n -
R_{n+1}}{\sqrt{R_n} + \sqrt{R_{n+1}}} =
\frac{a_n}{\sqrt{R_n} + \sqrt{R_{n+1}}} \geq
\frac{a_n}{2\sqrt{R_n}},
\]
故 $\sum_n \frac{a_n}{\sqrt{R_n}} \leq 2\sum_n(\sqrt{R_n} -
\sqrt{R_{n+1}}) = 2\sqrt{R_1} < \infty$（裂项相消）。然而
$\frac{a_n/\sqrt{R_n}}{a_n} = \frac1{\sqrt{R_n}} \to \infty$：
新级数收敛，却比原级数无穷地更大。
不存在最慢的收敛级数；任一固定族上的
比较判别法永远不可能\omterm{def:b2:metric:complete}{完备}。

\textbf{22.} 对递减正数列 $(a_n)$，将项
按相邻 $2$ 的幂之间分组：
\[
2^{k}a_{2^{k+1}} \leq \sum_{n=2^k}^{2^{k+1}-1} a_n \leq
2^ka_{2^k} .
\]
对 $k$ 求和：若 $\sum 2^ka_{2^k}$ 收敛，则
$\sum a_n$ 的部分和有界（从而收敛）；若
$\sum a_n$ 收敛，则 $\sum_k 2^{k+1}a_{2^{k+1}} \leq 2\sum_n a_n <
\infty$。对 $a_n = \frac1{n(\ln n)^\beta}$，$2^ka_{2^k} =
\frac{1}{(k\ln 2)^\beta}$，而 $\sum k^{-\beta}$ 收敛当且仅当
$\beta > 1$，这正是第 18 题的分界线，且无需
积分。

\textbf{23.} 由积分比较，
\[
\sum_{n > N}\frac{1}{n(\ln n)^2} \geq
\int_{N+1}^{\infty}\frac{\dd t}{t(\ln t)^2} =
\frac{1}{\ln(N+1)},
\]
当 $N = 10^6$ 时其值 $\approx 0.0724$：经过一百万项之后，
余项仍超过 $0.07$——级数收敛，但没有
任何直接求和能给出它的和。与
第 13 题对照：那里一次渐近修正，就从一百项
中换来七位数字：知道级数\emph{如何}收敛
比知道它收敛更有价值。

\textbf{24.} $\sum\frac1{n\ln n}$：发散（$\alpha = 1$，
$\beta = 1$，第 18 题）。$\sum\frac1{n^{1.01}}$：收敛
（黎曼，$\alpha > 1$）。$\sum\frac{(\ln n)^{100}}{n^{1.001}}$：
收敛（$\alpha = 1.001 > 1$，$\beta = -100$，第 16 题）。
$\sum\frac1{n\ln n(\ln\ln n)^3}$：收敛（第 19 题的
模式：原函数 $-\frac12(\ln\ln t)^{-2}$，有限
极限）。

\textbf{25.} 棣莫弗把 $\sin(2n{+}1)
\theta_k$ 为零这一事实转化为多项式在
$\cot^2\theta_k$ 处为零，而韦达定理读出了
仅凭分析只能估计的精确根和（第 1--5 题）。
夹逼需要两侧的\emph{精确}值 $\frac{n(2n-1)}3$ 与
$\frac{2n(n+1)}3$——等价量会循环论证，
因为全部要点正是常数
$\frac{\pi^2}6$（第 7--8 题）。第四部分全程依赖
\cref{ch:b2:comparison}中的级数与积分比较，
含对数的原函数承担分类的工作（第
18--19 题）。第 21 题粉碎了万能比较
判别法的梦想：每个收敛级数之下都还有另一个
无穷更慢的级数——伯特兰式的尺度会
越来越精细地刻画分界线，却永远无法到达。亮点有：欧拉的 $\zeta(2)
= \frac{\pi^2}6$，及其上一层 $\zeta(4) =
\frac{\pi^4}{90}$（第 8、10 题），以及伯特兰
分类（第 16--18 题）；$\zeta(2)$ 在
\cref{ch:b2:fourier}中再次出现，其中帕塞瓦尔恒等式
在一行之内由锯齿波的傅里叶级数重证它——一个
常数，两种文明。
\end{solution}
